%% ============================================================
%% Reframed Introduction — "When do supra-word tokens help?"
%% Draft companion to paper/reframe_outline.md.
%% Grounded in the English neural-BPB result: supra-word tokens buy
%% sequence compression at a small predictive-quality cost, and the
%% payoff depends on HOW they are selected (frequency vs MDL) and on
%% language type. Empirical figures marked [EN] are confirmed on
%% English; \TODO markers await fi/tr/zh.
%% ============================================================

\newcommand{\TODO}[1]{\textcolor{todocolor}{[\textbf{TODO}: #1]}}

\section{Introduction}
\label{sec:intro}

Almost every modern tokenizer refuses to place a token above the word.
Byte-Pair Encoding \citep[BPE;][]{sennrich-etal-2016-neural}, WordPiece
\citep{schuster1999japanese,wu2016googles}, and Unigram
\citep{kudo-2018-subword} all pre-tokenize on whitespace and then learn units
strictly \emph{within} the resulting word-like chunks.  Whitespace is a hard
floor: no matter how often \texttt{of the} or \texttt{New York} recurs, these
methods will never store it as a single token.  This convention is so
universal that it is rarely stated as a design choice at all — yet it is one,
and an under-motivated one.  It was inherited from early machine-translation
pipelines built for English, a language whose short, mostly monomorphemic,
space-delimited words make the word a convenient atom.  Whether the whitespace
floor is the \emph{right} place to stop has, until recently, gone essentially
untested.

Recent work has begun to question it.  SuperBPE \citep{liu2025superbpe} lets BPE
merges cross whitespace in a second training stage, producing multi-word
``superword'' tokens, and reports downstream gains (up to $+4\%$ across 30
tasks) with substantially less inference compute.  From a different tradition,
the LiB tokenizer \citep{yang2020less,yang2024rethinking} derives sub- and
supra-word units \emph{jointly} from a single Minimum Description Length (MDL)
objective, and \citet{yang2022eyefixations} show that its segmentations predict
human eye fixations during reading better than word-based ones — behavioural
evidence that multi-word units are psycholinguistically real
\citep{wray2002formulaic,arnon2021starting}.  Supra-word tokens are, suddenly,
no longer obviously a bad idea.

And yet the standard \emph{intrinsic} metrics of tokenization quality say they
are.  Under Description Length and under $n$-gram bits-per-character (BPC), both
LiB and SuperBPE rank \emph{worse} than the whitespace-bounded baselines they
extend.  This is the puzzle we take up: \textbf{if supra-word tokens help real
language models, why do intrinsic tokenization metrics penalize them?}

Our first contribution is to show that a large part of this penalty is a
\emph{measurement artefact}, not a property of the tokens.  Both Description
Length and $n$-gram BPC punish a tokenizer for having many rare types: DL adds
$\sim\!\log_2 N$ bits for every singleton type, and a Kneser–Ney $n$-gram model
has too little evidence to estimate a rare multi-word type, so its surprisal
explodes through data sparsity.  Supra-word tokens are, by construction, rarer
than the subwords they replace, so both metrics charge them for redundancy that
a better estimator would not see.  The tell is already visible in the data: the
same supra-word tokens \emph{help} a 2-gram model but \emph{hurt} a 3-gram model,
for both LiB and SuperBPE — a signature of the estimator, not the tokenization.

Our second contribution is to measure supra-word tokens on rulers that do not
have this bias.  We train a small neural language model on each tokenization and
report \textbf{bits-per-byte} (BPB) — cross-entropy over the test stream
normalized by raw UTF-8 bytes, the standard tokenization-fair compression
metric — alongside \textbf{fertility} (tokens per byte), which determines
sequence length and inference cost.  The neural metric removes the $n$-gram
sparsity pathology; the fertility metric exposes the efficiency dimension for
which supra-word tokens were designed.

The picture that emerges is a \emph{trade-off}, not a free lunch, and it is
worth stating plainly because it corrects the field in both directions.  On
English, supra-word tokens deliver a large, real compression benefit —
SuperBPE encodes the test text in the fewest tokens of any method
(\textasciitilde$0.18$ tokens/byte vs.\ $0.20$ for BPE) [EN] — but they do
\emph{not} improve predictive quality: LiB and SuperBPE both trail the
subword baselines on neural BPB, and turning supra-words on changes BPB only
marginally [EN].  In other words, the intrinsic metrics were right that
supra-word tokens cost something on the quality axis; they were wrong about
\emph{how much}, and they were silent about the compression benefit that
justifies the cost.  The honest case for supra-word tokens is an
efficiency-and-downstream case, not an intrinsic-quality case.

Our third contribution is to show that \emph{how} supra-word tokens are chosen
matters as much as \emph{whether} they are allowed, and that this is where a
principled objective earns its place.  SuperBPE selects supra-words by
frequency and compresses well; LiB selects them to minimize description length
and, on English, admits so many rare, low-yield multi-word types that it ends
up \emph{less} compressive than plain BPE despite spending 77\% of its
vocabulary on supra-words [EN].  The MDL objective is thus a double-edged tool:
it gives LiB a language-agnostic stopping criterion absent from every baseline
— on Turkish it halts at \TODO{29k} tokens because no further unit reduces
description length, rather than filling an arbitrary budget — but description
length is not the same objective as sequence compression, and optimizing it
does not automatically buy shorter sequences.  We use LiB and SuperBPE together
as two points in the space of supra-word selection rules to map when the
supra-word level pays off.

Finally, we ask \emph{for which languages} the trade-off is favourable.
The whitespace floor is not merely an English default; it interacts with
morphology.  In analytic languages the compressible redundancy lies mostly
\emph{above} the word (collocations), while in agglutinative and compounding
languages it lies \emph{within} it (morphemes) — so the same mechanism that
discovers ``\_of\_the'' in English should discover morpheme-length units in
Finnish or Turkish.  We evaluate across seven typologically diverse languages
— English, German, Spanish, Finnish, Turkish, Arabic, and Chinese — and find
that the fraction of supra-word tokens LiB discovers tracks the
analytic--synthetic continuum (\TODO{2\% for Chinese, 65--69\% for
agglutinative, 76--81\% for analytic and templatic}), giving a data-driven
typological diagnostic and localizing where a supra-word budget is best spent.

\paragraph{Contributions.}
\begin{itemize}
  \item We show that the intrinsic penalty against supra-word tokens under
        Description Length and $n$-gram BPC is largely a measurement artefact of
        rare-type sparsity, demonstrated by the 2-gram/3-gram reversal shared by
        LiB and SuperBPE.
  \item We provide a tokenization-fair neural evaluation (bits-per-byte plus
        fertility) that removes the artefact and reveals the true shape of the
        trade-off: supra-word tokens buy sequence compression at a small
        predictive-quality cost, not an intrinsic-quality gain.
  \item We show that the compression payoff depends on the supra-word
        \emph{selection rule}, contrasting frequency-based (SuperBPE) with
        MDL-based (LiB) admission, and identify over-admission of rare
        collocations as a failure mode of the MDL criterion.
  \item We present LiB as a principled MDL route to a joint sub-/supra-word
        vocabulary with an information-gain stopping criterion (Turkish
        saturation), released as a HuggingFace-compatible Rust implementation.
  \item We introduce the supra-word ratio as a typological diagnostic and
        characterize the trade-off across seven typologically diverse languages.
\end{itemize}
